\section{Related Works}
\label{sec:Related Works}
%-------------------------------------------------------------------------
\textbf{Continual Learning.}
Depending on the domain variations of incremental data, existing continual learning methods mainly focus on addressing \textit{i.e.,} Class Incremental Learning (CIL)~\cite{de2023effectiveness, jodelet2023class, belouadah2019il2m, yan2021dynamically, liu2023class} and Task Incremental Learning (TIL)~\cite{oren2021defense, mallya2018packnet, zheng2023preventing}. %Specifically, the CIL usually treats the subsets divided from single domain (\emph{e.g.,} one dataset) as incremental data, while the TIL usually directly uses data from multiple domains ((\emph{e.g.,} various datasets)). 
Existing efforts in this area have been made by developing various architectures~\cite{de2021continual}, including memory-based, regularization-based and dynamic-based models.  
% Memory-based
Memory-based methods~\cite{shin2017continual, lopez2017gradient, prabhu2020gdumb, lavda2018continual, ostapenko2022continual, rebuffi2017icarl, isele2018selective} retain the historical knowledge by storing them in a memory bank, which will be accessed and updated in incremental learning. However, the continuously increasing learned data usually poses a burden on the memory bank, resulting in limited lifelong learning ability. 
Regularization-based methods add explicit regularization terms on weights~\cite{aljundi2018memory, kirkpatrick2017overcoming, lee2017overcoming, zenke2017continual} or data~\cite{dhar2019learning, douillard2020podnet, hou2019learning, li2017learning} to balance between the older and new tasks. They are usually used as an auxiliary trick in memory-based or dynamic models to alleviate the forgetting issue.
Dynamic methods~\cite{douillard2022dytox, hu2023dense, ye2023self, yan2021dynamically, yoon2017lifelong, aljundi2017expert} address continual learning by incrementally adding new parameters on the baseline, such as neurons, branches or prediction heads. Dynamic methods usually perform favorably against the other two pipelines. However, like memory-based methods, the dynamic architecture often incurs large-scale model sizes, limiting the models' efficiency.
Despite the promising performance of the approaches aforementioned, addressing the crucial capability of AI agents, namely zero-shot transfer to unseen knowledge, remains challenging and complex to integrate into existing popular pipelines. In this paper, we propose incorporating the dynamic architecture on vision-language models to boost their memorization of historical knowledge and alleviate the degradation of zero-shot transfer ability.        
The highly related work is ZSCL~\cite{zheng2023preventing}, which uses parameter regularization in the continual learning of large-scale models. In contrast to the fully finetuning strategy in ZSCL, our method proposed an incremental MoE adapter to decrease the tuned parameters and enhance the collaboration of historically learned adapters and ongoing ones. 





%-------------------------------------------------------------------------

% attention->NLP
\begin{figure*}[t]
	\centering
	\includegraphics[width=0.98\linewidth]{figs/framework.pdf}
 \vspace{-4pt}
        \caption{Overall framework of the proposed method. (a) At the training stage, CLIP's image and text encoders $(\mathcal{F}_I,\mathcal{F}_T)$ take input samples from \textbf{Task \textit{t}}. 
        In each of transformer blocks, there is a MoE-Adapters, whose input is the tokens $ \textbf{\textit{x}}^t $ from MHSA. The router takes the task-specific \texttt{[CLS]} token $ \textbf{\textit{c}}^t $ as input and produces experts' weights $W_i^t$ and $W_j^t$ to combine the expert's output. DDAS is trained using only images via the MSE loss defined by Eq.~\ref{eq:chooseid1}. (b) At the inference stage, the proposed DDAS determines the data distribution by comparing the distribution $\{d^t\}_{t=1}^T$ in each autoencoder of the \textbf{task-agnostic} images. It can automatically assign the testing data into MoE-Adapters or original CLIP to predict with either seen or unseen data. 
        }
	% \caption{The framework of our method (a) the train stage,  (b) the inference stage, (c) the Auto-Choose module. During training (assuming \textbf{task 2} is ongoing), the images/text are sent to the image/text encoder respectively to train the MoE-Apdaters with task-dependent \textbf{router 2}. Furthermore, the image is sent to the Auto-chooser to train the \textbf{Autoencoder 2}. During inference, \textbf{task agnostic} images are fed into Auto-Choose to select router id, to assign the incoming testing data into learned MoE-adapters or pretrained CLIP.}
	\label{fig:overview}
 \vspace{-9pt}
\end{figure*}
\noindent\textbf{Parameter Efficient Fine-Tuning.}
%Recently, AI has stepped into the era of large-scale models. Built upon the general Transformer \cite{vaswani2017attention} and trained with billions of website data, large language models (LLM)~\cite{radford2019language, brown2020language, wei2022emergent, chen2021evaluating, ouyang2022training, touvron2023llama} or vision-language models (VLM)~\cite{radford2021learning, zhou2022learning, gao2023clip, liu2023visual, peng2023kosmos} show great potentiality to chat and interact like a ``real man''. 
In the realm of Natural Language Processing (NLP), fine-tuning large-scale models (\emph{e.g.,} 175B GPT-3~\cite{brown2020language}) imposes significant burdens in both parameter complexity and time consumption. Thus, several parameter-efficient fine-tuning methods~\cite{zhang2020side, hu2021lora, jia2022visual, houlsby2019parameter, karimi2021compacter} have been explored, which only set a few trainable parameters and fine-tune them for efficiency. Among these methods, LoRA \cite{hu2021lora} and Compacter~\cite{karimi2021compacter} reduce the number of trainable parameters by attaching low-rank hyper-complex adapter layers or sharing adapter parameters across layers, respectively.
The success of efficient tuning strategies in NLP promotes their applications on vision-language models~\cite{ju2022prompting, zhou2022conditional, gao2023clip, sung2022vl, zhang2021tip} like CLIP~\cite{gao2023clip}.
%In computer vision, the challenge of fine-tuning a large number of parameters persists in vision-language models such as CLIP.  ~\cite{liu2023class}
%A latest work~\cite{liu2023class} investigates different adaptation strategies (\emph{e.g.,} linear adapter, self-attention adapter) for vision-language model, which enables adaptation to new tasks in continual learning. 
Recently, Liu \textit{et al.}~\cite{liu2023class} introduce efficient tuning strategies in the continual learning of CLIP, which uses trainable adapters and a parameter retention strategy for downstream task adaptation and historical knowledge memorization, respectively.
However, this method is only applied to CIL and ignores the zero-shot ability of the original CLIP. 
In this paper, we propose a novel parameter-efficient tuning approach on CLIP to boost both the anti-forgetting and zero-shot abilities in continual learning. Our model can flexibly adapt to CIL and TIL and achieve promising performance even trained with few data (namely few-shot continual learning).
%Although prompt-based~\cite{ju2022prompting, zhou2022conditional} and adapter-based~\cite{gao2023clip, sung2022vl, zhang2021tip} methods are explored to seek a better trade-off between accuracy and efficiency, they are confined to handling individual tasks, posing challenges in their direct application to CL. 
%, we explored a parameter-efficient fine-tune method that enhances the performance of continual learning.



%-------------------------------------------------------------------------
\noindent\textbf{Mixture of Experts.}
The MoE~\cite{jacobs1991adaptive} contains multiple experts and a routing network. It aggregates the expert outputs via a weighted strategy by the routing network. Based on the sparse architecture of MoE~\cite{shazeer2017outrageously}, some methods~\cite{riquelme2021scaling, mustafa2022multimodal,fedus2022switch, daxberger2023mobile} are proposed to decrease computational costs and improve model capacity. This technique is also introduced to continual learning to mitigate the forgetting issue. For example, Aljundi \textit{et al.}~\cite{aljundi2017expert} propose to train multiple backbones as experts and automatically feed the test samples to a relevant expert. Chen \textit{et al.}~\cite{chen2023lifelong} utilize the pre-trained experts and gates to store previous knowledge. These methods have demonstrated MoE's promising performance in continual learning.
We propose an incremental MoE-Adapters for continual learning with CLIP. We use adapters as experts to increase the adaption speed and introduce an incremental expert interaction strategy to facilitate the collaboration of experts during continual learning. 






